\vspace{-0.5em}
\section{Related Work}
\vspace{-0.5em}
Extensive research has been done on applying LLMs to lengthy texts, with three main areas of focus: \textbf{Length Extrapolation}, \textbf{Context Window Extension}, and \textbf{Improving LLMs' Utilization of Long Text}. While seemingly related, it's worth noting that progress in one direction doesn't necessarily lead to progress in the other. For example, extending the context size of LLMs doesn't improve the model's performance beyond the context size, and neither approach ensures effective use of the long context. Our \method framework primarily lies in the first category, where LLMs are applied to text significantly exceeding the pre-training window size, potentially even of infinite length. We do not expand the attention window size of LLMs or enhance the model's memory and usage on long texts. The last two categories are orthogonal to our focus and could be integrated with our techniques.

\textbf{Length extrapolation} aims to enable language models trained on shorter texts to handle longer ones during testing. A predominant avenue of research targets the development of relative position encoding methods for Transformer models, enabling them to function beyond their training window. One such initiative is Rotary Position Embeddings (RoPE)~\citep{su2021roformer}, which transforms the queries and keys in every attention layer for relative position integration. Despite its promise, subsequent research~\citep{alibi, chen2023extending} indicated its underperformance on text that exceeds the training window. Another approach, ALiBi~\citep{alibi}, biases the query-key attention scores based on their distance, thereby introducing relative positional information. While this exhibited improved extrapolation, our tests on MPT models highlighted a breakdown when the text length was vastly greater than the training length. Current methodologies, however, have yet to achieve infinite length extrapolation, causing no existing LLMs to fit for streaming applications.

\textbf{Context Window Extension} centers on expanding the LLMs' context window, enabling the processing of more tokens in one forward pass. A primary line of work addresses the training efficiency problem. Given the attention to computation's quadratic complexity during training, developing a long-context LLM is both a computational and memory challenge. Solutions have ranged from system-focused optimizations like FlashAttention~\citep{dao2022flashattention, dao2023flashattention2}, which accelerates attention computation and reduces memory footprint, to approximate attention methods~\citep{zaheer2020bigbird,beltagy2020longformer,wang2020linformer,kitaev2020reformer} that trade model quality for efficiency. Recently, there has been a surge of work on extending pre-trained LLMs with RoPE~\citep{chen2023extending,kaiokendev,blocntkaware,peng2023yarn}, involving position interpolation and fine-tuning. However, all the aforementioned techniques only extend LLMs' context window to a limited extent, which falls short of our paper's primary concern of handling limitless inputs.

\textbf{Improving LLMs' Utilization of Long Text} optimizes LLMs to better capture and employ the content within the context rather than merely taking them as inputs. As highlighted by \citeauthor{liu2023lost} and \citeauthor{longchat2023}, success in the previously mentioned two directions does not necessarily translate to competent utilization of lengthy contexts. Addressing this effective usage of prolonged contexts within LLMs is still a challenge. Our work concentrates on stably harnessing the most recent tokens, enabling the seamless streaming application of LLMs.
